\par\clearpage\SourceLeaf{152}\ArithmeticHead{67}{R}{IVa. DAS AGM. NACH DEM 30. MAI 1799.}
\HistoricalFoot{}{9\textsuperscript{\(\ast\)}}
\indent Von diesem Beweise führt nun der folgende Gedankengang weiter. Durch
die Tatsache, dass die Reihe (\(\alpha\)) das absolute Glied der Entwicklung von
\((1+\sin^2\varphi)^{-\frac12}\) darstellt, wird \textsc{Gauss} auf die ihm aus der theoretischen Astronomie
bekannte Entwicklung von Ausdrücken der Form \((1+z\cos^2\varphi)^{-\frac12}\) geführt
(vergl. oben S. 26, 27), und wir finden auf S. 9 der Scheda (Werke X, 1,
S. 185) die Entwicklung von \((1-z\cos\varphi)^{-\frac12}\) mit der Darstellung der Koeffizienten
in der Form von unendlichen Reihen, die \textsc{Gauss} direkt aus \textsc{Eulers}
\textit{Calculus integralis} I, § 279, Probl. 33 (Opera, ser. I, vol. 11, S. 165) entnehmen
konnte, mit der daraus folgenden Feststellung, dass das absolute Glied der
Entwicklung von \((4+4z+z^2\cos^2\varphi)^{-\frac12}\) durch den reziproken Wert von \(M(1,1+z)\)
gegeben wird. — Um diesen Zusammenhang vertiefen zu können, fragt \textsc{Gauss}
zunächst wieder nach dem asymptotischen Verhalten von \(M(1,x)\) für grosse
Werte von \(x\). Wahrscheinlich hat er dieses Verhalten mit Hilfe des Ellipsenquadranten
erforscht, für den in der \textsc{Eulerschen} Abhandlung \textit{Animadversiones
etc.} (siehe oben S. 64) die Gleichung
\[
q=1+(\log2-\tfrac14)p^2-\tfrac12p^2\log p+(p^4)
\]
\noindent (Opuscula varii arg. II, § LIX, S. 163) gegeben wird, wo \((p^4)\) Glieder mit
der vierten und höheren Potenzen von \(p\) andeutet\footnote{Bestätigt wird die Darstellung des Textes durch die Tatsache, dass \textsc{Gauss} in seinem 1812 erworbenen Exemplar der Opuscula auf dem hintern Schutzblatte folgendes bemerkt hat:

»Formula art. LIX, p. 163 concinnius ita exhibetur:
\[
1+\log\frac4p\left\{\frac12p^2+\frac{1\,.\,1\,.\,3}{2\,.\,2\,.\,4}p^4+\cdots\right\}
=\frac1{\dfrac1{M(1,\sqrt{1-p^2})}+\left(\dfrac1{M(1,\sqrt{1-p^2})}-\dfrac{2q}\pi\right)\log\dfrac1x}\mathpunct{\vcenter{\hbox{,}}}
\]
ubi \(x\) determinatur per aequationem
\[
2x^{\frac12}+2x^{\frac92}+2x^{\frac{25}2}+\cdots=\sqrt{\frac p{M(1,\sqrt{1-p^2})}}\text{«}
\]
Wir haben nur der kürzern Schreibweise zuliebe für die von \textsc{Gauss} hingeschriebenen Reihenentwicklungen die Ausdrücke durch das agM. und den Ellipsenquadranten \(q\) gesetzt, wie sie in den obigen Gleichungen (1) bis (4) angegeben sind.}, die in Verbindung mit
der obigen Gleichung (3) sofort
\begingroup\makeatletter\tagsleft@true\let\veqno\leqno\makeatother
\begin{equation*}
\lim\limits_{p\to0}M(1,p)\log\frac4p=\frac12\pi\tag*{(8)}
\end{equation*}
\endgroup
\noindent oder das gleichwertige\par
\clearpage\SourceLeaf{153}\ArithmeticHead{68}{L}{L. SCHLESINGER, ÜBER GAUSS’ FUNKTIONENTHEORETISCHE ARBEITEN.}
\begingroup\makeatletter\tagsleft@true\let\veqno\leqno\makeatother
\begin{equation*}
\lim\limits_{x\to\infty}\frac1xM(1,x)\log4x=\frac12\pi,\tag*{(8a)}
\end{equation*}
\noindent liefert (vergl. die Bemerkungen Werke X, 1, S. 268—270). Diese Formel
kommt zwar in der Scheda Ac nicht explizite vor, sie liegt aber den in dieser
von S. 10 ab (Werke X, 1, S. 186) aufgezeichneten Formeln zu Grunde. Wenn
man nämlich durch die Gleichung (siehe (6), S. 65)
\[
\nu=4z-16z^3+56z^5+\cdots
\]
\noindent die uniformisierende Variable \(z\) einführt, so folgt für \(\nu=\frac1x\) aus (8a)
\begin{equation*}
\lim\limits_{x\to\infty}\frac1xM(1,x)\log\frac1z=\frac12\pi\tag*{(9)}
\end{equation*}
\noindent was zu dem Ansatz (Werke X, 1, S. 186, Gl. (7))
\begin{equation*}
M(1,x)=\frac{\dfrac\pi2\left(x-\dfrac\alpha x-\dfrac\beta{x^3}-\cdots\right)}{\log\dfrac1z}\tag*{(10)}
\end{equation*}
\noindent hinleitet, von dem es in der schon in Helmstedt am 14. Dezember 1799 geschriebenen
Tagebuchnotiz 101. heisst: »Dass das agM. als der Quotient zweier
transzendenter Funktionen darstellbar sei, hatten wir bereits früher gefunden«.
Wenn diese Tagebuchnotiz nun fortfährt: »nunmehr haben wir entdeckt, dass
die eine dieser Funktionen [nämlich der Zähler] auf Integralgrössen zurückgeführt
werden kann«, so ist damit der Inhalt der S. 11 der Scheda (Werke
X, 1, S. 187, Art. 4) gemeint. \textsc{Gauss} bezeichnet daselbst den reziproken Wert
des Zählers in (10) mit \(Q\), entwickelt diese Grösse nach fallenden Potenzen
von \(x\) und erkennt in ihr den von \(\varphi\) unabhängigen Teil von \((x^2-\cos^2\varphi)^{-\frac12}\),
den er nunmehr mit dem Werte des Integrals
\begin{equation*}
\frac2\pi\int_0^1\frac{dr}{\sqrt{1-r^2\,.\,x^2-r^2}}\tag*{(11)}
\end{equation*}
\endgroup
\noindent identifiziert\footnote{Man vergl. hierzu die Darstellung in der Abhandlung von 1800, abgedruckt Werke III, S. 370—371, die aus etwas späterer Zeit stammt. Die Berechnung der Koeffizienten \(\alpha,\beta,\ldots\) des Zählers in (10) siehe bei \textsc{Geppert}, a. a. O.\textsuperscript{83)}, S. 108.}.

Nun erweist sich aber die für \(Q\) gefundene Reihenentwicklung als mit der
für \(\frac\nu{M(1,\sqrt{1-\nu^2})}\) übereinstimmend, so dass sich also\par
\clearpage\SourceLeaf{154}\ArithmeticHead{69}{R}{IVa. DAS AGM. NACH DEM 30. MAI 1799.}
\begingroup\makeatletter\tagsleft@true\let\veqno\leqno\makeatother
\begin{equation*}
\frac\nu Q=M(1,\sqrt{1-\nu^2}),\quad\nu=\frac1x\mathpunct{\vcenter{\hbox{,}}}\tag*{(12)}
\end{equation*}
\noindent und folglich
\begin{equation*}
\frac1{M(1,\sqrt{1-\nu^2})}=\frac2\pi\int_0^1\frac{dr}{\sqrt{(1-r^2)(1-\nu^2r^2)}}\tag*{(13)}
\end{equation*}
\noindent ergibt (Werke X, 1, S. 187, Gl. (20)), ein Resultat, das \textsc{Gauss} später (1818,
\textit{Determinatio attractionis}, Art. 16, Werke III, S. 331) in viel direkterer Weise
bewiesen hat. — Dieser Zusammenhang des agM. mit dem vollständigen elliptischen
Integral erster Gattung spielt in den weiteren Untersuchungen von
\textsc{Gauss} eine sehr bedeutsame Rolle; ob sich aber die Tagebuchnotiz 102. vom
23. Dezember 1799: »das agM. ist selbst eine Integralgrösse« auf diese Feststellung
bezieht, ist darum zweifelhaft, weil ihr Wortlaut fast genau mit der
Bemerkung in der Abhandlung von 1800 (Werke III, S. 370) übereinstimmt,
wo es heisst: »auf diese Weise sind also unsere agMittel auf Integralgrössen
zurückgeführt«, und diese Bemerkung an die lineare Differentialgleichung zweiter
Ordnung
\begin{equation*}
x(x^2-1)y''+(3x^2-1)y'+xy=0\tag*{(14)}
\end{equation*}
\noindent anknüpft, der die reziproken Werte von \(M(1,x)\) und \(M(1,\sqrt{1-x^2})\) Genüge
leisten. — Diese Differentialgleichung wird in den Werke X, 1, S. 181—183
abgedruckten Zetteln aufgestellt, die (vergl. die Bemerkungen Werke X, 1,
S. 273, 274) sicher in diesen Tagen der Jahreswende 1799 zu 1800 geschrieben
sind.

Die Gleichung (10) ergibt in Verbindung mit (12) die wichtige Beziehung
(Werke X, 1, S. 190, Art. 7)
\begin{equation*}
\log z=-\tfrac12\pi\frac{M(1,\sqrt{1-\nu^2})}{M(1,\nu)}\mathpunct{\vcenter{\hbox{,}}}\tag*{(15)}
\end{equation*}
\noindent die die uniformisierende Variable \(z\) unmittelbar durch die agM. darstellt,
während diese Variable in den Leistenotizen (siehe oben Gl. (5)) gewissermassen
nur implizite definiert war. — Die Gleichung (15) in Verbindung mit
den an derselben Stelle aufgezeichneten Gleichungen
\begin{equation*}
\left\{\begin{aligned}
\frac\nu{M(1,\sqrt{1-\nu^2})}&=2(2z+2z^9+\cdots)(1+2z^4+2z^{16}+\cdots),\\
\frac1{M(1,\sqrt{1-\nu^2})}&=(2z+2z^9+\cdots)^2+(1+2z^4+2z^{16}+\cdots)^2
\end{aligned}\right.\tag*{(16)}
\end{equation*}
\endgroup\par
\clearpage\SourceLeaf{155}\ArithmeticHead{70}{L}{L. SCHLESINGER, ÜBER GAUSS’ FUNKTIONENTHEORETISCHE ARBEITEN.}
\noindent führen zu einem gewissen Abschluss der Theorie des agM.\begingroup\renewcommand{\thefootnote}{90a)}\footnote{Zufolge der Beziehung \(2p(x^2)r(x^2)=r(x)^2\) (vergl. den folgenden Abschnitt IVb) stimmt die erste der Gleichungen (16) mit (5) überein.}\endgroup\addtocounter{footnote}{-1}.
— Wann \textsc{Gauss}
diese Aufzeichnungen, die bis zur Seite 15 der Scheda Ac gehen (Werke X, 1,
S. 190—193), gemacht hat, lässt sich nicht ganz genau feststellen. Da sich aber
auf S. 22 der Scheda die Werke II, S. 311 abgedruckte Aufzeichnung über
ternäre Formen findet\footnote{Vergl. \textsc{Bachmann}, Abh. 1 dieses Bandes, S. 26.}, und \textsc{Gauss} auf diesen Gegenstand bezügliche Untersuchungen
in der Tagebuchnotiz 103. vom 13. Februar 1800 erwähnt, so kann
man wohl annehmen, dass er die durch die Gleichungen (15), (16) bezeichneten
Entdeckungen vor dem 13. Februar 1800 gemacht hat, und dass dann
eine Pause in den Arbeiten über elliptische Funktionen eingetreten ist.

Offenbar ist \textsc{Gauss} erst durch die Entdeckung des Zusammenhanges zwischen
dem agM. und dem vollständigen elliptischen Integral erster Gattung
auf die Bedeutung des Integrals erster Gattung für die allgemeine Theorie
der elliptischen Funktionen aufmerksam geworden. Vorher war er wohl der
Ansicht, dass das Analogon der Kreisfunktionen nur bei denjenigen Integralen
zu suchen sei, die ähnlich wie das Kreisbogenintegral die Rektifikation
gewisser Kurven liefern, also wie das Integral \(\displaystyle\int\frac{dx}{\sqrt{1-x^4}}\) die der Lemniskate, das
allgemeine Integral zweiter Gattung die der Ellipse. Aber die weiteren Entwicklungen
aus der Theorie des agM. lieferten ihm noch weit mehr als diese
blosse Anregung zur Beschäftigung mit dem Integral erster Gattung.

Die Einsicht, dass die 1798 für die lemniskatische Periode gefundenen
Reihen (4) des Abschn. III (S. 39) sich aus den jetzt für das allgemeine agM.
gefundenen Reihen (16) für \(\nu=\sqrt{\tfrac12}\) ergeben, enthüllte ihm die Bedeutung des
auf den linken Seiten von (16) auftretenden Ausdrucks \(\frac\pi{M(1,\sqrt{1-\nu^2})}\) als der
einen Periode (um in der uns geläufigen Terminologie zu sprechen) des allgemeinen
elliptischen Integrals erster Gattung.

Grosse Schwierigkeiten mögen ihm aber die Realitätsverhältnisse bereitet
haben. Schon für die Wahl des Moduls bot das lemniskatische Integral keinen
rechten Anhaltspunkt, weil für dieses der Modul bald gleich \(\sqrt{-1}\), also imaginär,
bald gleich \(\sqrt{\tfrac12}\), also reell und kleiner als Eins erscheint\raisebox{.45ex}{.} Wir sehen
in der Tat, wie \textsc{Gauss} zwischen den verschiedenen Normalformen\par
\clearpage\SourceLeaf{156}\ArithmeticHead{71}{R}{IVa. DAS AGM. NACH DEM 30. MAI 1799.}
\begin{center}\small
\begin{tabular}{@{}ccc@{}}
\(\displaystyle\int\frac{du}{\sqrt{1+\mu\cos u}}\mathpunct{\vcenter{\hbox{,}}}\)&
\(\displaystyle\int\frac{du}{\sqrt{1+\mu^2\sin^2u}}\mathpunct{\vcenter{\hbox{,}}}\)&
\(\displaystyle\int\frac{dx}{\sqrt{(1-\alpha x^2)(1-\beta x^2)}}\mathpunct{\vcenter{\hbox{,}}}\)\\
(Scheda Ac, S. 26&(Scheda Ac, Titel&(\textit{Tagebuch} Nr. 105 und\\
Werke X, 1, S. 194)&Werke X, 1, S. 184)&Werke VIII, S. 96)
\end{tabular}\end{center}
\noindent lange Zeit hin- und herschwankt. Wie mühevoll mag ihm erst die Erkenntnis
dessen geworden sein, dass der Quotient der beiden Perioden eines elliptischen
Integrals erster Gattung keine reelle Grösse sein kann!

Einen Einblick in diese Zeit des Langens und Bangens gewährt die folgende
Aufzeichnung, die sich auf einem Zettel ohne Datum findet, die aber
offenbar an den Anfang des Jahres 1800 zu setzen ist\footnote{Am 6. Mai 1800 hat \textsc{Gauss} nämlich schon völlige Klarheit erreicht, vergl. \textit{Tagebuch} 105 und weiter unten.}:

»Der Radikalfehler, woran meine bisherigen Bestrebungen, den Geist der
elliptischen Funktion zu verkörpern, gescheitert sind, scheint \textls[100]{der} zu sein,
dass ich dem Integral
\[
\int\frac{d\varphi}{\sqrt{(1-e^2\sin^2\varphi)}}
\]
\noindent die Bedeutung als Ausdruck eines endlichen Teils der Kugelfläche habe unterlegen
wollen, während es wahrscheinlich nur einen unendlich schmalen Kugelsektor
ausdrücken soll.«

Was \textsc{Gauss} hier als »wahrscheinlich« bezeichnet, wäre also, dass jenes
Integral (für \(e\) reell und kleiner als Eins) zwei \textls[100]{reelle} Perioden mit inkommensurablem
Verhältnis besitzt\footnote{Bekanntlich hat auch \textsc{Abel} eine zeitlang mit diesem Paradoxon gekämpft; vergl. \textsc{Gundelfinger}, Sitzungsberichte der Berl. Akad. 1898, S. 344.}; aber wie ein Blitz leuchtet dieser Irrtum
in die Gedankenwerkstatt von \textsc{Gauss} hinein: wir sehen, dass er die komplexe
Variable auf der Kugel interpretiert und den dem Periodenparallelogramm
entsprechenden Fundamentalbereich entwirft, wir sehen, dass er aus der Existenz
zweier reeller Perioden auf das Vorhandensein einer unendlich kleinen
Periode schliesst. Und alles das im Jahre 1800!

Um Ostern 1800 ist \textsc{Gauss} seiner ursprünglichen Absicht gemäss nach
Braunschweig zurückgekehrt; Ostersonntag fiel auf den 13. April und die
Tagebuchnotiz 104 vom 27. April ist die erste, die wieder in Braunschweig
geschrieben ist. Bald darauf hat er nun auch alle Schwierigkeiten überwunden,
und wir sehen ihn vom 6. Mai ab wieder von Entdeckung zu Entdeckung\par
\clearpage\SourceLeaf{157}\ArithmeticHead{72}{L}{L. SCHLESINGER, ÜBER GAUSS’ FUNKTIONENTHEORETISCHE ARBEITEN.}
\noindent schreiten, mit denen er in den wenigen Wochen bis gegen Mitte Juni 1800
den gewaltigen Bau seiner Theorie der elliptischen und der Modulfunktionen
errichtet.
